%; whizzy chapter -dvi
% -initex iniptex -latex platex -format platex -bibtex jbibtex -fmt fmt
% 以上 whizzytex を使用する場合の設定。

%     Tokyo Debian Meeting resources
%     Copyright (C) 2012 Junichi Uekawa
%     Copyright (C) 2011, 2015 Nobuhiro Iwamatsu

%     This program is free software; you can redistribute it and/or modify
%     it under the terms of the GNU General Public License as published by
%     the Free Software Foundation; either version 2 of the License, or
%     (at your option) any later version.

%     This program is distributed in the hope that it will be useful,
%     but WITHOUT ANY WARRANTY; without even the implied warranty of
%     MERCHANTABILITY or FITNESS FOR A PARTICULAR PURPOSE.  See the
%     GNU General Public License for more details.

%     You should have received a copy of the GNU General Public License
%     along with this program; if not, write to the Free Software
%     Foundation, Inc., 51 Franklin St, Fifth Floor, Boston, MA  02110-1301 USA

%  preview (shell-command (concat "evince " (replace-regexp-in-string "tex$" "pdf"(buffer-file-name)) "&"))

%%ここからヘッダ開始。

\documentclass[mingoth,a4paper]{jsarticle}
\usepackage{monthlyreport}
% 日付を定義する、毎月変わります。
\newcommand{\debmtgyear}{2018}
\newcommand{\debmtgmonth}{3}
\newcommand{\debmtgdate}{24}
% started from zero:
% (let ((year 2013) (month 7)) (+ (* (- year 2005) 12) month -1))
\newcommand{\debmtgnumber}{161}

% Needed to import pandoc-generated LaTeX documents.
% See https://stackoverflow.com/questions/40438037/tightlist-error-using-pandoc-with-markdown
\providecommand{\tightlist}{%
  \setlength{\itemsep}{0pt}\setlength{\parskip}{0pt}}

% tikz picture の為のマクロ設定
\usepackage[dvipdfmx]{graphicx}
\usepackage{tikz}

\begin{document}

\begin{titlepage}
\thispagestyle{empty}
% タイトルページ:編集必要な部分は最初のマクロに飛ばすこと

\vspace*{-2cm}
第\debmtgnumber{}回 東京エリア Debian 勉強会資料\\
\hspace*{-2cm}
\includegraphics{image2012-natsu/dotdeb.pdf}\\
\hfill{}\debmtgyear{}年\debmtgmonth{}月\debmtgdate{}日

% ここはアップデートすること
% 全角文字にしないとフォントのサイズが合わないので注意
\rotatebox{10}{\fontsize{30}{30} {\gt　GO／機械学習 特集}}\\

\vspace*{-2cm}
\hfill{}\includegraphics[height=6cm]{image200502/openlogo-nd.eps}
\end{titlepage}

\newpage

\begin{minipage}[b]{0.2\hsize}
 \definecolor{titleback}{gray}{0.9}
 \colorbox{titleback}{\rotatebox{90}{\fontsize{80}{80} {\gt デビアン勉強会} }}
\end{minipage}
\begin{minipage}[b]{0.8\hsize}
\hrule
\vspace{2mm}
\hrule
\begin{multicols}{2}
\tableofcontents
\end{multicols}
\vspace{2mm}
\hrule
\end{minipage}

\dancersection{最近のDebian関連のミーティング報告}{杉本 典充}

\subsection{第159回東京エリアDebian勉強会}

2018年1月20日(土)に第159回東京エリアDebian勉強会を開催しました。会場は東銀座にある朝日ネットさんをお借りして行いました。参加者は3名でした。

各自が用意した課題に対してハックを行いました。


\subsection{OSC 2018 Tokyo/Spring}

2018年2月24日(土)に明星大学様でOSC 2018 Tokyo/Springが開催されました。東京エリアDebian勉強会はこのイベントに出展し、セミナー発表とブース展示を行いました。
イベント全体の来場者は、2日間合計で1200名と主催者から発表されています。


セミナー発表は、表題「apacheにdjangoのwebアプリでデプロイしてみよう」として杉本が発表し、聴講者は21名でした。


ブース展示では、来場者へのDebian利用の実態に関するヒアリング、広報活動を行いました。


\dancersection{事前課題}{杉本 典充}

今回の事前課題は以下です。

\begin{enumerate}
\item Hack Timeは何をしますか
\item 機械学習の環境に触れたことがある方は、利用した環境を教えてください。
\end{enumerate}

この課題に対して提出いただいた内容は以下です。

\begin{multicols}{2}
{\small
\begin{prework}{ hiromiso }
  \begin{enumerate}
  \item $B7h$a$F$$$^$;$s(B
  \item scikit-learn$B!"(BSVN$B!"(BC/C++$B$G$N;0AX9=B$%Q!<%;%W%H%m%s<BAu!"(Btf-idf$B$rMQ$$$?<+A38@8l%/%i%9%?%j%s%0(B
  \end{enumerate}
\end{prework}

\begin{prework}{ yy\_y\_ja\_jp }
  \begin{enumerate}
  \item Salsa$BBP1~!$%Q%C%1!<%899?7(B
  \item ($B2sEz$J$7(B)
  \end{enumerate}
\end{prework}

\begin{prework}{ ysaito }
  \begin{enumerate}
  \item go$B=q$-$^$9(B
  \item ($B2sEz$J$7(B)
  \end{enumerate}
\end{prework}

\begin{prework}{ henrich }
  \begin{enumerate}
  \item debootstrap$B$N%P%0:F8=4D6-:n@.(B
  \item ($B2sEz$J$7(B)
  \end{enumerate}
\end{prework}

\begin{prework}{ ichinomoto }
  \begin{enumerate}
  \item dm200$B$G(Bsystemd$B<~$j$N@_Dj@0M}$J$I(B
  \item ($B2sEz$J$7(B)
  \end{enumerate}
\end{prework}

\begin{prework}{ John Paul Adrian Glaubitz John Paul Adrian Glaubitz }
  \begin{enumerate}
  \item Hack on Debian Ports stuff, like always :P.
  \item no answer.
  \end{enumerate}
\end{prework}

\begin{prework}{ Roger Shimizu }
  \begin{enumerate}
  \item maintain packages
  \item $BJ9$$$?$3$H$,$"$k$@$1!#(Bcaffe$B$J$I(B
  \end{enumerate}
\end{prework}

\begin{prework}{ dictoss }
  \begin{enumerate}
  \item salsa$B$X(Bgit$B%j%]%8%H%j$r0\$9(B
  \item $B4D6-$K?($l$?$3$H$O$"$j$^$;$s(B
  \end{enumerate}
\end{prework}

\begin{prework}{ su\_do }
  \begin{enumerate}
  \item $B@^3Q$J$N$G(BDebiab$B4D6->e$G$N5!3#3X=,4D6-$N9=C[Ey=PMh$l$P$HB8$8$^$9!#(B
  \item ($B2sEz$J$7(B)
  \end{enumerate}
\end{prework}

}
\end{multicols}

%\dancersection{Debian Trivia Quiz}{username}
%
%Debianの昨今の話題についてのQuizです。
%
%今回の出題範囲は\url{debian-devel-announce@lists.debian.org} や \url{debian-news@lists.debian.org}などに投稿された内容からです。
%
%\begin{multicols}{2}
%%; whizzy-master ../debianmeetingresume201211.tex
% $B0J>e$N@_Dj$r$7$F$$$k$?$a!"$3$N%U%!%$%k$G(B M-x whizzytex $B$9$k$H!"(Bwhizzytex$B$,MxMQ$G$-$^$9!#(B
%

\santaku
{DebConf13 $B$N3+:ECO$H3+:EF|$O!)(B}
{$BF|K\!"El5~ET(B 6$B7n(B20$BF|(B}
{$B%K%+%i%0%"(B $B%^%J%0%"(B 7$B7n(B8-14$BF|(B}
{$B%9%$%9!"%t%)!<%^%k%-%e(B 8$B7n(B11-18$BF|(B}
{3}
{$B%K%+%i%0%"$O(BDebConf12$B$N3+:ECO$G$9!#(B
DebConf13$B$O%9%$%9$N%-%c%s%WCO$G3+:E$G$9!#(B
6/20$B$O3'$5$sM=Dj$r6u$1$F$*$-$^$7$g$&!#(B}

\santaku
{$B@$3&$N(BWeb$B%5!<%P$G:G$b?M5$$N$"$k(BLinux $B%G%#%9%H%j%S%e!<%7%g%s(B(W3Techs$BD4$Y(B)$B$O!)(B}
{CentOS}
{Debian}
{Ubuntu}
{B}
{\url{http://w3techs.com/technologies/history_details/os-linux}$B$K7k2L$N%0%i%U$,$"$j$^$9!#(B
$B8=:_(B Linux $B$r;HMQ$7$F$$$k(B web $B%5!<%P$N(B 32.9\% $B$,(B Debian $B$rMxMQ$7$F$*$j!"$=$N3d9g$O8=:_$bA}2C$rB3$1$F$$$k$=$&$G$9!#(B}

\santaku
{Debian $B%+!<%M%k%A!<%`$N%a%s%P!<$G$"$j!"(Bkernel.org $B$N(B 3.2.y $B0BDjHG7ONs$N%a%s%F%J$G$b$"$k(B Ben Hutchings $B$5$s$,<!4|(B Debian $B0BDjHG$H0l=o$K=P2Y$5$l$k(B Linux $B%+!<%M%k$K(B (3.2 $B7ONs$N(B mainline $B$K$OL5$$(B) $BDI2C5!G=$,Ek:\$5$l$kM=Dj$G$"$k$H=R$Y$F$$$^$9!#(B
$BB?$/$NDI2CE@$NCf$K4^$^$l$J$$$b$N$O2?!)(B}
{PREEMPT\_RT}
{Hyper-V guest drivers$B$N6/2=(B}
{ARM64/AArch64$B%"!<%-%F%/%A%c%5%]!<%H(B}
{C}
{Hyper-V guest drivers$B$O(Bmainline kernel$B$G(B3.2$B$K$b4^$^$l$F$$$^$9$,!"$h$j2~A1$5$l$?(B3.4$B$+$i$N=$@5$,F3F~$5$l$^$9!#(B
PREEMPT\_RT$B$O%O!<%I%j%"%k%?%$%`$r<B8=$9$k$?$a$N(BPatch$B!"(B
linux-image-rt-amd64 , linux-image-rt-686-pae $B$N(Bmetapackage$B$G;HMQ$G$-$^$9!#(B
$B?7$7$$(BARM 64$B%S%C%H%"!<%-%F%/%A%c%5%]!<%H$O(Bmainline kernel 3.7$B$+$i(B}

\santaku
{Wookey$B$5$s$,%"%J%&%s%9$7$?(Balpha$BHG$N(BDebian port arm64 image$B$O!)(B}
{Debian/Ubuntu port image}
{Debian/KFreeBSD port image}
{Debian/GnuHurd port image}
{A}
{self-bootstrapp(non x86)$BBP1~$H$N$3$H$G$9!#(B\url{http://wiki.debian.org/Arm64Port}$B$G%9%F!<%?%9$,3NG'$G$-$^$9!#(B}

\santaku
{700,000$BHVL\$N%P%0$,Js9p$5$l$?F|$rEv$F$k(B700000thBugContest$B$N7k2L$,=P$^$7$?!#$=$NM=A[F|$HJs9pF|$O!)(B}
{2012/12/12$B$rM=A[$7$?(BDavidPrevot}
{$BM=A[F|(B:2013/02/04$B!"Js9pF|(B:2013/02/14}
{$BM=A[F|(B:2013/02/07$B!"Js9pF|(B:2013/02/14}
{$BM=A[F|(B:2013/02/14$B!"Js9pF|(B:2013/02/07}
{C}
{$B:G$b6a$$(B2013/02/14$B$rM=A[$7$?(BChristian Perrier$B$5$s$,Ev$F$^$7$?!#7k2L$O(B\url{http://wiki.debian.org/700000thBugContest}$B$G8x3+$5$l$F$$$^$9!#(B
$B$^$?!"(B800,000/1,000,000$BHVL\$N%P%0$,Js9p$5$l$kF|$rEv$F$k%3%s%F%9%H(B\url{http://wiki.debian.org/800000thBugContest}$B$b3+:E$5$l$F$$$^$9!#(B}

\santaku
{master.debian.org$B$,?7$7$$5!3#$K0\9T$5$l$^$7$?!#$3$l$O2?$N%5!<%P$G$7$g$&$+(B $B!)(B}
{@debian.org$B$N%a!<%k%5!<%P(B}
{$B%Q%C%1!<%8$N%^%9%?!<%5!<%P(B}
{$B%Q%C%1!<%8$N%9%]%s%5!<(B(mentor)$B$rC5$9%5!<%P(B}
{A}
{$B8E$$%5!<%P$O%G%#%9%/>c32Ey$,$"$C$?$N$G!"<wL?$HH=CG$5$l!"%G!<%?$,B;<:$9$kA0$K?7$7$$%5!<%P$K0\9T$5$l$^$7$?!#(Bftp-master.debian.org$B$O(BDebian$B$N(B official package $B%j%]%8%H%j$G$9!#%Q%C%1!<%8$N%9%]%s%5!<(B(mentor)$B$rC5$9$N$O(Bmentors.debian.net$B!#(B }

\santaku
{pbuilder$B$K(Bclang support$B$,DI2C$5$l$^$7$?!#C/$,=q$$$?%Q%C%A$G$7$g$&$+!)(B}
{Sylvestre Ledru}
{Junichi Uekawa}
{Hideki Yamane}
{C}
{Debian$B$N(BClang$B%5%]!<%H$OCe!9$H?J$s$G$$$^$9!#(B}

\santaku
{DPN - 2013$BG/(B3$B7n(B4$BF|9f$K<h$j>e$2$i$l$?F|K\$N%$%Y%s%H$O(B}
{Open Source Conference 2013 Tokyo/Spring}
{Open Source Conference 2013 Hamamatu}
{Open Source Conference 2013 Tokushima}
{A}
{\url{http://henrich-on-debian.blogspot.jp/2013/02/open-source-conference-2013-tokyospring.html} $B>\:Y$O8e$[$I!#(B}


%\end{multicols}


% % (query-replace-regexp "<.*?>" "")
% % (query-replace-regexp "^[    ]\+" "")

%-------------------------------------------------------------------------------
\dancersection{go / debian での機械学習環境構築について}{ysaito}
%-------------------------------------------------------------------------------

\subsection{はじめに}
機械学習という分野では, python, あるいは, R が人気を二分する.
しかし, 第三の選択肢として, Go による機械学習環境がある.
Go で機械学習を行うメリットは, 静的型安全や, goroutine, channel, などの並行処理のメリットがあるが, 最も注目すべき点は, インフラ, システムプログラミングに対する親和性であると考える.

今回は, 機械学習そのものの解説には踏み込まず, Go/Debian による機械学習環境の構築に触れたい.

環境は debian 9.4 (stretch) を利用する.

\subsection{書籍の紹介}
Machine Learning With Go (Packt publishing)

\url{https://www.packtpub.com/big-data-and-business-intelligence/machine-learning-go}

\subsection{kuberetes, docker の準備}
VMサポートは, 有効にせずローカルでの実行を前提として minikube を導入する.

docker リポジトリの導入
\begin{commandline}
apt install curl apt-transport-https ca-certificates software-properties-common
curl -fsSL https://download.docker.com/linux/debian/gpg | sudo apt-key add -
\end{commandline}

fingerprint の確認
\begin{commandline}
apt-key fingerprint 0EBFCD88

pub   rsa4096 2017-02-22 [SCEA]
      9DC8 5822 9FC7 DD38 854A  E2D8 8D81 803C 0EBF CD88
uid           [ unknown] Docker Release (CE deb) <docker@docker.com>
sub   4096R/F273FCD8 2017-02-22 [S]
\end{commandline}

docker の導入
\begin{commandline}
apt install docker-ce
docker version
Docker version 18.03.0-ce, build 0520e24

# 入力データを処理するアルゴリズムがのった docker image を導入する
docker pull dwhitena/goregtrain:single
\end{commandline}


minikube, kubectl の導入
\begin{commandline}
curl -Lo minikube https://storage.googleapis.com/minikube/releases/latest/minikube-linux-amd64 && \
chmod +x minikube && sudo mv minikube /usr/local/bin/
curl -Lo kubectl \
https://storage.googleapis.com/kubernetes-relesase/release/\
$(curl -s https://storage.googleapis.com/kubernetes-release/release/stable.txt)/bin/linux/amd64/kubectl && \
chmod +x kubectl && mv kubectl /usr/local/bin/

minikube version
minikube version v0.25.2

kubectl version
Client Version: version.Info{Major:"1", Minor:"9", GitVersion:"v1.9.6",
GitCommit:"9f8ebd171479bec0ada837d7ee641dec2f8c6dd1", GitTreeState:"clean",
BuildDate:"2018-03-21T15:21:50Z", GoVersion:"go1.9.3", Compiler:"gc",
Platform:"linux/amd64"}
\end{commandline}

minikube 用の設定と minikube の開始
\begin{commandline}
export MINIKUBE_WANTUPDATENOTIFICATION=false
export MINIKUBE_WANTREPORTERRORPROMPT=false
export MINIKUBE_HOME=$HOME
export CHANGE_MINIKUBE_NONE_USER=true
mkdir $HOME/.kube || true
touch $HOME/.kube/config

export KUBECONFIG=$HOME/.kube/config

minikube start --vm-driver=none

\end{commandline}

pachyderm の導入
\begin{commandline}
curl -o /tmp/pachctl.deb -L \
https://github.com/pachyderm/pachyderm/releases/download/v1.7.0rc4/pachctl_1.7.0rc4_amd64.deb && \
sudo dpkg -i /tmp/pachctl.deb

pachctl deploy local
\end{commandline}

\subsection{pachyderm とは}
pachydermとは, データのバージョン管理や, 機械学習の処理をパイプラインでつなぐことができる go 製のソフトウェアである.
詳しくは, \url{https://pachyderm.io} にて情報がある.

\subsection{構成する概要}

diabetes.csv から, go による機械学習アルゴリズムにより、モデル（csvから得られた関数のパラメータ）を生成し, model.json
新規に得られた入力, 1.json から予測, 1.json を出力する

\begin{commandline}
     +-------------------------------------------------------------+
     |                                                             |
     |                                                             |
     |  +----------+    +----------+    +------------+             |
     |  |          |    |          |    |            |             |
     |  |  学   習 |  　| モ デ ル |    |  モ  デ  ル|             |
     |  |          |    |          |    |            |             |
+-----> |          +--> |          |--> |            |             |
     |  diabetes.csv    | ./goregtrain  | model.json |             |
     |  |          |    |          |    |            |             |
     |  +----------+    +----------+    +-----+------+             |
     |                                        |                    |
     |                  +----------+    +-----v------+  +--------+ |
     |                  |          |    |            |  |        | |
     |                  | attributes    | 予  測     |  | 予  測 | |
 +--------------------> |          +--> |            +> |        | |
     |                  |          |    |            |  |        | |
     |                  | 1.json   |    | ./goregpredict| 1.json | |
     |                  +----------+    +------------+  +--------+ |
     |                                                             |
     +-------------------------------------------------------------+
\end{commandline}
     
\subsection{入力に使用するデータ}
\url{https://github.com/PacktPublishing/Machine-Learning-With-Go/tree/master/Chapter09/building_a_scalable_pipeline/example2}

\subsection{go によって pachyderm へ接続しリポジトリを作る}
\begin{commandline}
// localhost 上の Kubernetes クラスタの pachyderm へ接続する
// デフォルトの pachyderm のポート 30650
c, err := client.NewFromAddress("0.0.0.0:30650")
if err != nil {
log.Fatal(err)
}
defer c.Close()
// 学習用のリポジトリを作成 "training."
if err := c.CreateRepo("training"); err != nil {
log.Fatal(err)
}
// 予測の入力用のリポジトリを作成する "attributes."
if err := c.CreateRepo("attributes"); err != nil {
log.Fatal(err)
}
// ２つのリポジトリに sanity check をおこなう
repos, err := c.ListRepo(nil)
if err != nil {
log.Fatal(err)
}
// リポジトリの数を確認する
if len(repos) != 2 {
log.Fatal("Unexpected number of data repositories")
}
\end{commandline}

コンパイルと実行
\begin{commandline}
go build
./create
pachctl list-repo
NAME CREATED SIZE
attributes 3 seconds ago 0B
training 3 seconds ago 0B
\end{commandline}

\subsection{attribute リポジトリ, training リポジトリへデータをセットする}

\begin{commandline}
// Pachyderm へ接続する
c, err := client.NewFromAddress("0.0.0.0:30650")
if err != nil {
log.Fatal(err)
}
defer c.Close()
// "attributes" データリポジトリにデータを "master" ブランチにコミットする処理を始める
commit, err := c.StartCommit("attributes", "master")
if err != nil {
log.Fatal(err)
}
// attributes に入れる JSON を開く
f, err := os.Open("1.json")
if err != nil {
log.Fatal(err)
}
// attributes へファイルをプットする
if _, err := c.PutFile("attributes", commit.ID, "1.json", f); err != nil {
log.Fatal(err)
}
// コミットを完了させる.
if err := c.FinishCommit("attributes", commit.ID); err != nil {
log.Fatal(err)
}
// "training" データリポジトリの "master" ブランチへデータをコミットする処理を始める.
commit, err = c.StartCommit("training", "master")
if err != nil {
log.Fatal(err)
}
// 学習用のデータセットファイルを開く
f, err = os.Open("diabetes.csv")
if err != nil {
log.Fatal(err)
}
// training データセットを 学習用データリポジトリに展開する.
if _, err := c.PutFile("training", commit.ID, "diabetes.csv", f); err !=
nil {
log.Fatal(err)
}
// コミットを完了させる.
if err := c.FinishCommit("training", commit.ID); err != nil {
log.Fatal(err)
}
\end{commandline}

コンパイルして実行する
\begin{commandline}
# 上記のコードをコンパイルする
go build
# 実行
./a

# リポジトリを表示する
pachctl list-repo
NAME CREATED SIZE
training 13 minutes ago 73.74KiB
attributes 13 minutes ago 210B

# training リポジトリの master ブランチのファイルを表示する
pachctl list-file training master
NAME TYPE SIZE
diabetes.csv file 73.74KiB

# attributes リポジトリの master ブランチのファイルを表示する
pachctl list-file attributes master
NAME TYPE SIZE
1.json file 210B
\end{commandline}

\subsection{model ステージのJSON を準備する}

\begin{commandline}
{
  "pipeline": {
    "name": "model"
  },
  "transform": {
    "image": "dwhitena/goregtrain:single",
    "cmd": [
      "/goregtrain",
      "-inDir=/pfs/training",
      "-outDir=/pfs/out"
    ]
  },
  "parallelism_spec": {
    "constant": "1"
  },
  "input": {
    "atom": {
      "repo": "training",
      "glob": "/"
    }
  }
}
\end{commandline}

\subsection{model ステージのJSON を準備する}

1.Pachyderm データパイプラインが model という名前であることを教える
2.Pachyderm へ, モデル作成に使うアルゴリズムを教える, docker イメージになっている線形回帰モデルを使用する( dwhitena/goregtrain:single ), 
3. そして, 学習用データファイルとのパイプラインを教える

\begin{commandline}
# パイプラインの作成
pachctl create-pipeline -f model.json

# pods の状態確認
kubectl get pods
NAME READY STATUS RESTARTS AGE
etcd-2142892294-38ptw 1/1 Running 0 2h
pachd-776177201-04l6w 1/1 Running 0 2h
pipeline-model-v1-p0lnf 2/2 Running 0 1m

# pachyderm 上の job 確認
pachctl list-job
ID OUTPUT COMMIT STARTED DURATION RESTART PROGRESS DL UL STATE
14f052ae-878d-44c9-a1f9-ab0cf6d45227 model/a2c7b7dfb44a40e79318c2de30c7a0c8
3 minutes ago Less than a second 0 1 + 0 / 1 73.74KiB 160B success

# データリポジトリを確認
pachctl list-repo
NAME CREATED SIZE
model 3 minutes ago 160B
training About an hour ago 73.74KiB
attributes About an hour ago 210B

# model master ブランチにあるファイルを確認
pachctl list-file model master
NAME TYPE SIZE k8s
model.json file 160B

# model.json の中身を確認する
pachctl get-file model master model.json
{
  "intercept": 152.13348416289818,
  "coefficients": [
    {
      "name": "bmi",
      "coefficient": 949.4352603839862
    }
  ]
}
\end{commandline}


\subsection{予測用のパイプラインを設定し結果を得る}

\begin{commandline}
{
  "pipeline": {
    "name": "prediction"
  },
  "transform": {
    "image": "dwhitena/goregpredict",
    "cmd": [
      "/goregpredict",
      "-inModelDir=/pfs/model",
      "-inVarDir=/pfs/attributes",
      "-outDir=/pfs/out"
    ]
  },
  "parallelism_spec": {
    "constant": "1"
  },
  "input": {
    "cross": [
    {
      "atom": {
        "repo": "attributes",
        "glob": "/*"
      }
    },
      {
        "atom": {
          "repo": "model",
          "glob": "/"
        }
      }
    ]
  }
}
\end{commandline}

\begin{commandline}
# prediction.json へのデータパイプラインを作成する
pachctl create-pipeline -f prediction.json

# pachyderm 上の job のリストを確認する
pachctl list-job
ID OUTPUT COMMIT STARTED DURATION RESTART PROGRESS DL UL STATE
03f36398-89db-4de4-ad3d-7346d56883c0
prediction/5ce47c9e788d4893ae00c7ee6b1e8431 About a minute ago Less than a
second 0 1 + 0 / 1 370B 266B success
14f052ae-878d-44c9-a1f9-ab0cf6d45227 model/a2c7b7dfb44a40e79318c2de30c7a0c8
19 minutes ago Less than a second 0 1 + 0 / 1 73.74KiB 160B success

# リポジトリの確認をする
pachctl list-repo
NAME CREATED SIZE
prediction About a minute ago 266B
model 19 minutes ago 160B
training About an hour ago 73.74KiB
attributes About an hour ago 210B

# prediction リポジトリの master ブランチの中身を確認する
pachctl list-file prediction master
NAME TYPE SIZE
1.json file 266B

# 1.json ファイルの中身を確認する
pachctl get-file prediction master 1.json
\end{commandline}

%
% 冊子にするために、4の倍数にする必要がある。
% そのための調整
\dancersection{メモ}{}
\mbox{}\newpage


\vspace*{15cm}
\hrule
\vspace{2mm}
\includegraphics[width=2cm]{image200502/openlogo-nd.eps}
\noindent \Large \bf Debian 勉強会資料\\
\noindent \normalfont \debmtgyear{}年\debmtgmonth{}月\debmtgdate{}日 \hspace{5mm}  初版第1刷発行\\
\noindent \normalfont 東京エリア Debian 勉強会 （編集・印刷・発行）\\
\hrule

\end{document}
